Se quieren construir señales reales mediante una constante y sólo los armónicos de índices \(2\) y \(3\) respecto de un periodo \(2\pi\). Los dos coeficientes del primer armónico deben ser cero. Se prescribe
\[
\frac1{2\pi}\int_0^{2\pi}f(t)\,dt=2,
\qquad
\frac1{2\pi}\int_0^{2\pi}f(t)^2\,dt=Q.
\]
La segunda cantidad es la media del cuadrado, no su raíz. La señal debe ser no constante y tener \(2\pi\) como menor periodo positivo.

Determine para qué valores de \(Q\) es posible la construcción. Describa las amplitudes y fases que quedan libres bajo esas condiciones, y exhiba una señal concreta para \(Q=5\). Justifique tanto los momentos prescritos como el periodo mínimo de la señal construida; no basta que \(2\pi\) sea uno de sus periodos.
